\section{Đồng hồ nói khác phép đếm}
\label{sec:ch08-dong-ho}

\index{đồng hồ!vì sao phải đo riêng}%
Mục trước đếm được chính xác tới từng FLOP. Mục này hỏi cái đồng hồ, và hai câu
trả lời không khớp nhau. Chỗ chúng lệch là nội dung của mục này và của mục sau.

Trước hết một chuyện về cách đo, vì một bảng thời gian không nói rõ cách đo thì
không lặp lại được. Mọi số dưới đây là trung vị sau một lần chạy làm nóng: chín
lần lặp ở các bảng chỉ đo forward, năm lần ở bảng đo cả forward lẫn backward, vì
bảng ấy đắt hơn một bậc. Cột \code{spread} bằng lớn nhất trừ nhỏ nhất, để một
hàng bị nhiễu tự lộ ra. Máy là một CPU 20 nhân và \code{torch.get\_num\_threads()} trả
về 20. Con số cuối ấy không phải chi tiết vặt, và cuối mục có một bảng cho thấy
vì sao.

\begin{measured}{một tầng con, forward, batch 8, 8 head}
\begin{minted}{text}
d    n     attention   spread      repo loop   nn.LSTM     attn/lstm
128  32    0.000470    0.000258    0.003149    0.000454    1.03
128  64    0.000470    0.000391    0.006256    0.000776    0.61
128  128   0.000786    0.000177    0.006428    0.001368    0.57
128  256   0.001860    0.000418    0.026182    0.002508    0.74
128  512   0.023154    0.014496    0.026852    0.004983    4.65
128  1024  0.087120    0.010140    0.068895    0.009825    8.87
512  32    0.000979    0.000335    0.004106    0.001723    0.57
512  64    0.002140    0.000327    0.008026    0.004157    0.51
512  128   0.003691    0.000508    0.022771    0.007345    0.50
512  256   0.009894    0.001401    0.036816    0.015056    0.66
512  512   0.035190    0.003143    0.086326    0.028690    1.23
512  1024  0.127809    0.010557    0.182810    0.057187    2.23
\end{minted}

Đo bằng \code{experiments/ch08\_clock.py} tại tag \code{ch08}.
\end{measured}

\subsection*{Chỗ đổi dấu không đi theo bề rộng}

\index{đồng hồ!chỗ đổi dấu không phụ thuộc d}%
Đọc cột \code{attn/lstm}, chỗ nó vượt qua 1 là chỗ attention bắt đầu đắt hơn.
Bỏ qua hàng \(n = 32\), nơi cả hai bên chỉ tốn vài phần mười nghìn giây và phần
chi phí cố định lấn át hết phần số học; ở đó con số 1.03 không nói gì về kiến
trúc cả.

Ở \(d = 512\), cột ấy đi từ 0.66 ở \(n = 256\) lên 1.23 ở \(n = 512\). Ở
\(d = 128\), nó đi từ 0.74 ở \(n = 256\) lên 4.65 ở \(n = 512\). Hai bề rộng
cách nhau bốn lần, và chỗ đổi dấu nằm trong \emph{cùng một khoảng},
giữa 256 và 512.

Đó không phải cái \eqref{eq:ch08-attention-vs-lstm} dự đoán. Theo FLOP, chỗ đổi
dấu ở \(n = 2d\), tức 256 ở bề rộng nhỏ và 1024 ở bề rộng lớn: nó phải \emph{đi
theo} bề rộng. Đo thì nó đứng yên.

Ở \(d = 512\) phép đếm nói attention còn rẻ tới 1024 và đồng hồ nói nó hết rẻ ở
quãng 512, tức đồng hồ khắc nghiệt hơn phép đếm đúng hai lần. Ở \(d = 128\) thì
ngược lại, phép đếm nói 256 và đồng hồ cho attention đi thêm được một quãng.

Một đại lượng đứng yên khi \(d\) đổi bốn lần thì không phải hàm của \(d\).
Mục~\ref{sec:ch08-cho-vo} là chỗ đi tìm xem nó là hàm của cái gì, và câu trả lời
giải thích luôn cả cột \code{spread}, vốn phình lên đúng ở hàng
\(d = 128, n = 512\): 0.014496 trên một trung vị 0.023154, tức gần hai phần ba.

\subsection*{Chữ \enquote{recurrent} phải nói rõ là cài đặt nào}

\index{đồng hồ!vòng lặp Python so với kernel gộp}%
Cột \code{repo loop} là encoder của chính repo này, một vòng lặp Python qua từng
\tn{bước thời gian}{time step}, và cột \code{nn.LSTM} là cùng phép tính ấy sau
một kernel đã gộp
sẵn. Ở forward, khoảng cách giữa hai cột còn đi theo bề rộng: quanh năm tới mười
lần ở \(d = 128\) và chỉ quanh hai tới ba lần ở \(d = 512\).

Ở forward cộng backward thì khoảng cách ấy đổi hẳn bậc, và đó là phép đo đáng
giá hơn, vì huấn luyện trả cả hai chiều:

\begin{measured}{forward cộng backward, \(d = 128\), batch 8}
\begin{minted}{text}
n     attention   repo loop   nn.LSTM     loop/lstm  attn/lstm
128   0.002371    0.036614    0.008458    4.3        0.28
256   0.012425    0.082913    0.015498    5.3        0.80
512   0.049383    0.856043    0.033927    25.2       1.46
\end{minted}

Đo bằng \code{experiments/ch08\_clock.py} tại tag \code{ch08}.
\end{measured}

Ở \(n = 512\), vòng lặp của repo chậm hơn kernel gộp 25.2 lần. Lý do thì nằm ở
autograd chứ không ở số học: vòng lặp Python dựng một nút đồ thị cho mỗi bước
thời gian, nên chiều ngược phải đi qua 512 nút thay vì một.

Chuyện ấy có hệ quả cho cách đọc cả chương. Nếu tôi lấy cột \code{repo loop} làm
\enquote{một tầng hồi quy} thì attention trông tốt hơn nhiều so với thực tế, và
kết luận của chương sẽ là một phát biểu về Python chứ không phải về kiến trúc.
Bảng in cả hai cột vì thế, và mọi so sánh trong chương này lấy cột
\code{nn.LSTM}.

Nó cũng đặt giá lên một lựa chọn của repo. Mọi tầng hồi quy trong repo này là
một vòng lặp Python viết tay, không phải \code{torch.nn.LSTM}, vì một cuốn sách
dựng lại từ đầu rồi lặng lẽ thay bằng thư viện ở đúng chương cần tốc độ thì đã
thôi dựng lại từ đầu. Cột \code{nn.LSTM} ở đây là để so, không phải để thay. Con
số 25.2 là cái giá của lựa chọn ấy, đo được, ở một chỗ.

\subsection*{Số head: món nợ của chương trước, trả kèm một đính chính}

\index{multi-head!đồng hồ không đồng ý với FLOP}%
Mục~\ref{sec:ch08-dem} chứng minh số head triệt tiêu khỏi phép đếm FLOP, và
\code{tests/test\_ch08.py} khẳng định điều đó thành một đẳng thức đúng. Số tham
số cũng không đổi theo \(h\), và chương trước đã có một test cho chuyện ấy.

Nên bảng dưới đây là một bảng mà cả hai cột kiểm được đều hằng, và chỉ đồng hồ
đổi:

\begin{measured}{quét số head ở \(\dmodel = 512\), FLOP và tham số đều hằng}
\begin{minted}{text}
every row holds 1,048,576 parameters

h    d_k   n=128       vs h=1   n=512       spread      vs h=1
1    512   0.003287    1.00     0.017986    0.004141    1.00
2    256   0.003602    1.10     0.020214    0.001773    1.12
4    128   0.003221    0.98     0.026352    0.002604    1.47
8    64    0.003421    1.04     0.035445    0.002492    1.97
16   32    0.003639    1.11     0.058174    0.010010    3.23
32   16    0.004835    1.47     0.096588    0.011568    5.37
\end{minted}

Đo bằng \code{experiments/ch08\_clock.py} tại tag \code{ch08}.
\end{measured}

Ở \(n = 512\), ba mươi hai head chậm hơn một head 5.37 lần. Cùng số phép tính,
cùng số tham số, cùng phần cứng.

Lý do thì không sâu: chia một phép nhân ma trận lớn thành \(h\) phép nhỏ hơn thì
mỗi phép có ít việc hơn để che chi phí không phải số học, và ở \(h = 32\) thì
\(\dk = 16\), tức các ma trận đủ hẹp để phần dựng và dọn lấn át. Câu của bài,
\enquote{the total computational cost is similar}, đúng theo phép đếm và bài nói
rõ nó đang đếm.

\index{multi-head!đính chính con số của chương 7}%
\textbf{Một đính chính cho chương trước.} Mục~\ref{sec:ch07-nhieu-dau} viết rằng
một phiên đo riêng thấy \(h = 16\) chậm hơn \(h = 1\) khoảng hai mươi phần trăm
ở \(\dmodel = 512\) trên CPU. Con số ấy sai. Đo tử tế thì là 3.23 lần ở
\(n = 512\) và 1.11 lần ở \(n = 128\), không phải hai mươi phần trăm ở bất kỳ độ
dài nào tôi thử, và bốn lần chạy độc lập trên máy rảnh cho 3.23, 3.03, 3.09 và
3.32 nên nó không phải nhiễu. Chương trước đã tự nói rằng nó không in con số ấy thành một
phát biểu chung vì nó phụ thuộc cài đặt và phần cứng, và sự thận trọng ấy đúng
chỗ. Cái sai là vẫn in nó ra.

Cột \code{n=128} thì không đọc được thành xu hướng: 1.00, 1.10, 0.98, 1.04,
1.11, 1.47 không đơn điệu, vì ở đó mỗi phép đo chỉ vài phần
nghìn giây và phần chi phí cố định lấn át. Tôi in nó ra để nói rằng hiệu ứng phụ
thuộc độ dài, và để không ai đọc cột \code{n=512} thành một hằng số của kiến
trúc.

\subsection*{Đổi số luồng thì đổi cả kết luận}

\index{đồng hồ!số luồng đổi chỗ giao nhau}%
Một bảng cuối, ngắn, và nó là lý do mọi bảng trên đều ghi số luồng.

\begin{measured}{cùng phép đo ở 20 luồng và ở 1 luồng, \(d = 512\)}
\begin{minted}{text}
threads  n     attention   nn.LSTM     attn/lstm
20       128   0.003935    0.007834    0.50
20       512   0.035344    0.031290    1.13
1        128   0.022032    0.041252    0.53
1        512   0.118236    0.141877    0.83
\end{minted}

Đo bằng \code{experiments/ch08\_clock.py} tại tag \code{ch08}.
\end{measured}

Ở \(n = 512\), hai mươi luồng cho 1.13 và một luồng cho 0.83. Tức ở cùng một
kiến trúc, cùng một độ dài, cùng một máy, attention đắt hơn hay rẻ hơn
\tn{mạng hồi quy}{recurrent network} \emph{tùy vào số luồng}. Đổi một tham số không có trong bảng 1, không có
trong \eqref{eq:ch08-flop-tang}, và không có trong bất kỳ phát biểu nào của bài,
thì kết luận đổi dấu.

Đó là lời cảnh báo tôi muốn để lại trước khi sang mục sau. Một bảng
\tn{độ phức tạp}{complexity}
nói về thuật toán. Một bảng thời gian nói về một thuật toán chạy trên một cỗ máy
cụ thể ở một cấu hình cụ thể, nên khi bạn gặp một bảng thời gian không ghi cỗ
máy và cấu hình ra, nó không nói về gì cả.
